% !TeX root = 20261028_exp91.tex
% !LW recipe=lablog XeLaTeX
\input{_support/style.tex}
\begin{document}
\reporttitle{2026-10-28}{EXP91: 두 TabPFN 설계안의 코드 구현과 동작 검증}
\section{목적과 실제 수행 범위}
Tabular Prior-data Fitted Network (TabPFN) 기반 두 모델의 학습·추론·모델 저장 및 복원 경로가 설계대로 동작하는지 검증한다. 사용자 요청은 BETA·CURE를 기존 모델 설계 보고서에 추가하고 1·2안 코드를 작성하는 것이었다. 승인된 연구 버전 v7을 유지하여 EXP91 구현 검증으로 등록했다. 이 실행은 Intrusion Detection System (IDS; 침입 탐지 시스템)의 CIC2018·ToN-IoT 성능 실험이나 문헌 모델과의 비교 결과가 아니다.

설계·문헌·수식·11개 그림은 \texttt{docs/\allowbreak{}\allowbreak{}lab\allowbreak{}log/\allowbreak{}\allowbreak{}repo\allowbreak{}rts/\allowbreak{}\allowbreak{}2026\allowbreak{}1028\_\allowbreak{}\allowbreak{}mod\allowbreak{}el\_\allowbreak{}\allowbreak{}desig\allowbreak{}ns.\allowbreak{}tex}에 정리했다. 본 문서는 실제 실행 상태와 실패 수정 근거를 별도로 기록한다. 코드 진입점은 \texttt{scripts/\allowbreak{}\allowbreak{}v7/\allowbreak{}\allowbreak{}exp91\allowbreak{}/\allowbreak{}\allowbreak{}run\_\allowbreak{}\allowbreak{}des\allowbreak{}igns.\allowbreak{}py}, 설정은 \texttt{configs/\allowbreak{}\allowbreak{}experime\allowbreak{}nts/\allowbreak{}\allowbreak{}exp9\allowbreak{}1\_\allowbreak{}\allowbreak{}protoc\allowbreak{}ol.\allowbreak{}json}이다.

\section{자료·환경·고정 설정}
합성자료는 seed 43으로 만든 3개 클래스다. 각 클래스 참라벨 6개로 총 18개, 비라벨 학습 36개, test 24개이며 모두 6차원 연속 특징이다. 행 ID와 동일 특징 벡터 group이 세 집합 사이에서 겹치지 않는지 확인했다. 비라벨 정답은 모델에 제공하지 않고 test 정답은 예측 후 채점에만 사용했다.

실제 TabPFN-v3 checkpoint와 NVIDIA RTX 4090을 사용했다. smoke는 estimator 1개, 학생 은닉 32·표현 16차원, 준비 1 epoch와 교대 갱신 1회·학습 1 epoch, 2안 최대 전문가 1개·2-fold로 줄였다. 이는 기본 학습 예산보다 작은 기능 검증용 설정이다. Python·패키지 버전과 checkpoint hash는 각 실행의 config.json에 기록했다.

\section{검증 결과}
단위 검증 10개가 모두 통과했다. 검사에는 soft 배정 목적함수의 일차 최적조건, 비균등 prior 반영, 수치·범주·결측 전처리, 참라벨/\allowbreak{}비라벨 특징 중복 거부, 동일 특징 group의 Out-of-fold (OOF) 분리, n=\allowbreak{}2의 정확한 Global fallback, 두 logit 보정의 상한·다른 logit 불변, 전문가 채택 경로, 학생의 가중치 갱신과 저장 복원이 포함됐다. 테스트용 분류기는 알고리즘 경로를 의도적으로 만들기 위한 대역이며 분류 성능의 근거로 쓰지 않는다.

실제 PFN 검증에서는 1안 encoder의 가중치 hash가 바뀌었고 교사 가중치는 유지됐다. 2안은 쌍 [1,2] 전문가 하나를 실제로 채택하여 Global과 전문가 결합 경로를 실행했다. 모든 PFN의 가중치 불변을 확인했으며 같은 8개 query batch에서 저장·복원 전후 최대 확률 차이는 두 모델 모두 0이었다. 학생은 교사와 복원 decoder 없이 배포할 수 있고, 2안은 native PFN fit state와 구조 JSON을 저장했다.

\begingroup
\footnotesize
\begin{longtable}{>{\raggedright\arraybackslash}p{\dimexpr 0.31373\linewidth-2\tabcolsep\relax}>{\raggedright\arraybackslash}p{\dimexpr 0.29412\linewidth-2\tabcolsep\relax}>{\raggedright\arraybackslash}p{\dimexpr 0.39216\linewidth-2\tabcolsep\relax}}
\toprule
항목 & 확인 결과 & 해석 범위 \\ \midrule
\endfirsthead
항목 & 확인 결과 & 해석 범위 \\ \midrule
\endhead
{}단위 검증 & 10/\allowbreak{}10 통과 & 수식·정보 경계·예외 경로 \\
{}1안 실제 PFN 학습·추론 & 완료, 학생 가중치 갱신 & 합성자료 파이프라인 \\
{}2안 실제 전문가 채택 & 1개 채택 & 순차 보정 경로가 실제 실행됨 \\
{}PFN 가중치 & 변경 없음 & context/\allowbreak{}cache 준비와 가중치 학습 구분 \\
{}같은 batch의 저장/\allowbreak{}복원 & 두 모델 확률 차이 0 & 모델 export 정상 \\
{}IDS 정확도·비용 우위 & 미검증 & 전체 데이터 실험 필요 \\
\bottomrule
\end{longtable}
\endgroup
\section{초기 실패와 수정}
첫 실행 \texttt{20261028\allowbreak{}\_\allowbreak{}\allowbreak{}194707\_\allowbreak{}\allowbreak{}two\_\allowbreak{}\allowbreak{}desi\allowbreak{}gns\_\allowbreak{}\allowbreak{}real\allowbreak{}\_\allowbreak{}\allowbreak{}pfn\_\allowbreak{}\allowbreak{}smo\allowbreak{}ke\_\allowbreak{}\allowbreak{}1e824\allowbreak{}bfb}는 모델 예측까지 수행했지만 2안의 저장·복원 검사에서 실패했다. 전체 24행 예측의 앞 8행과 복원 후 8행 batch를 비교했으며 최대 확률 차이가 약 0.000313이었다. 같은 모델에서도 query batch 크기가 달라질 때 발생한 차이와 저장 오류를 구분하지 못한 검사 조건이었다. 허용 오차를 확대하지 않고 저장 전후 모두 같은 8행 batch를 사용하도록 고쳤으며, batch 크기에 따른 차이는 별도 진단으로 보존했다.

수정 실행 \texttt{20261028\allowbreak{}\_\allowbreak{}\allowbreak{}194827\_\allowbreak{}\allowbreak{}two\_\allowbreak{}\allowbreak{}desi\allowbreak{}gns\_\allowbreak{}\allowbreak{}real\allowbreak{}\_\allowbreak{}\allowbreak{}pfn\_\allowbreak{}\allowbreak{}smo\allowbreak{}ke\_\allowbreak{}\allowbreak{}fixed\allowbreak{}\_\allowbreak{}\allowbreak{}probe\_\allowbreak{}\allowbreak{}0\allowbreak{}0ac57e9}는 완료됐다. 이어 CLI에서 만든 학생 export를 package API에서도 읽을 수 있도록 전처리기의 Python 객체 대신 기본 상태 딕셔너리를 저장하도록 바꿨다. 최종 실행 \texttt{20261028\allowbreak{}\_\allowbreak{}\allowbreak{}195409\_\allowbreak{}\allowbreak{}two\_\allowbreak{}\allowbreak{}desi\allowbreak{}gns\_\allowbreak{}\allowbreak{}port\allowbreak{}able\_\allowbreak{}\allowbreak{}exp\allowbreak{}ort\_\allowbreak{}\allowbreak{}smok\allowbreak{}e\_\allowbreak{}\allowbreak{}d0ff03\allowbreak{}79}에서 동일 검증을 완료했고 독립 package import로 두 모델을 불러와 예측했다. 이전 실패·중간 성공의 소스 snapshot과 실행 기록은 보존했다.

\section{실행과 후속 검증}
일반 실행은 NPZ의 X\_\allowbreak{}l·y\_\allowbreak{}l·X\_\allowbreak{}u·X\_\allowbreak{}test를 받으며 y\_\allowbreak{}test와 행 ID는 선택 항목이다. 범주형 열은 --categorical-columns로 지정하고 학생·쌍 설정은 --config JSON의 student와 pair 항목으로 조정한다. 실제 실험은 등록 실행을 통해 입력 hash·라벨 수·비라벨 수·코드 snapshot·예측·모델·비용을 함께 저장한다. 실행 명령과 기본값은 모델 설계 보고서의 구현 절에 정리했다.

두 코드의 다음 검증은 동일 라벨과 비라벨 접근 예산에서 단독 PFN·단순 증류·대표점 대조·단순 쌍 전문가 및 문헌 방법과 비교하는 것이다. 특히 n=\allowbreak{}2에서 2안의 보정이 꺼지는 비율, 학생의 희귀 공격 pseudo-label 소실, 정상 오탐과 공격별 precision·recall을 확인한다. 합성자료의 F1과 수초 단위 실행 시간을 IDS 우위나 실제 서비스 비용으로 해석하지 않는다. BETA·CURE 자체의 재현은 이번 구현에 포함되지 않았다.

출처: \texttt{docs/\allowbreak{}\allowbreak{}lab\allowbreak{}log/\allowbreak{}\allowbreak{}evid\allowbreak{}ence/\allowbreak{}\allowbreak{}202\allowbreak{}61028/\allowbreak{}\allowbreak{}ex\allowbreak{}p91\_\allowbreak{}\allowbreak{}vali\allowbreak{}dation.\allowbreak{}j\allowbreak{}son} 출처: \texttt{docs/\allowbreak{}\allowbreak{}lab\allowbreak{}log/\allowbreak{}\allowbreak{}evid\allowbreak{}ence/\allowbreak{}\allowbreak{}202\allowbreak{}61028/\allowbreak{}\allowbreak{}ex\allowbreak{}p91\_\allowbreak{}\allowbreak{}unit\allowbreak{}\_\allowbreak{}\allowbreak{}tests.\allowbreak{}l\allowbreak{}og} 출처: \texttt{tabpfn/\allowbreak{}\allowbreak{}r\allowbreak{}esults/\allowbreak{}\allowbreak{}v\allowbreak{}7/\allowbreak{}\allowbreak{}exp91/\allowbreak{}\allowbreak{}20261028\allowbreak{}\_\allowbreak{}\allowbreak{}195409\_\allowbreak{}\allowbreak{}two\_\allowbreak{}\allowbreak{}desi\allowbreak{}gns\_\allowbreak{}\allowbreak{}port\allowbreak{}able\_\allowbreak{}\allowbreak{}exp\allowbreak{}ort\_\allowbreak{}\allowbreak{}smok\allowbreak{}e\_\allowbreak{}\allowbreak{}d0ff03\allowbreak{}79/\allowbreak{}\allowbreak{}confi\allowbreak{}g.\allowbreak{}json} 출처: \texttt{tabpfn/\allowbreak{}\allowbreak{}r\allowbreak{}esults/\allowbreak{}\allowbreak{}v\allowbreak{}7/\allowbreak{}\allowbreak{}exp91/\allowbreak{}\allowbreak{}20261028\allowbreak{}\_\allowbreak{}\allowbreak{}195409\_\allowbreak{}\allowbreak{}two\_\allowbreak{}\allowbreak{}desi\allowbreak{}gns\_\allowbreak{}\allowbreak{}port\allowbreak{}able\_\allowbreak{}\allowbreak{}exp\allowbreak{}ort\_\allowbreak{}\allowbreak{}smok\allowbreak{}e\_\allowbreak{}\allowbreak{}d0ff03\allowbreak{}79/\allowbreak{}\allowbreak{}summa\allowbreak{}ry.\allowbreak{}json}


\end{document}
